\chapter{Boundary Handling and Safe Gradients}
\label{ch:bf-boundary-gradients}

Boundary handling is part of the posterior target.  Parameters can leave the
stationary, positive-definite, or economically admissible region during HMC
warmup.  A robust implementation must handle those proposals without crashing
the chain or returning non-finite gradients.

\section{Finite Invalid Returns}
\label{sec:bf-finite-invalid-returns}

The preferred policy is:
\begin{enumerate}[label=(\roman*)]
  \item detect invalid regions before fragile factorizations where possible;
  \item return either the exact support-rejection value ($-\infty$) or a declared finite fallback scalar for a modified target;
  \item if a finite fallback scalar is used, return the exact derivative of that same scalar;
  \item record the invalid reason in diagnostics;
  \item keep the behavior stable under JIT compilation.
\end{enumerate}

This policy is not a mathematical trick to make an invalid model valid. The exact support-rejection contract and the finite modified-target contract are different mathematical objects and must be named as such. A finite fallback may keep sampler control flow in a compiled finite-valued regime, but it does not by itself guarantee rejection or preserve the exact support-restricted target.

\section{Boundary Checklist}
\label{sec:bf-boundary-checklist}

Every HMC target should state its policy for prior support, parameter
transforms, stationarity or determinacy restrictions, covariance pivots, solve
residuals, non-finite likelihoods, and non-finite gradients.
